\apendice{Anexo de sostenibilización curricular}

\section{Introducción}

Este anexo incluye una reflexión personal sobre los aspectos de sostenibilidad abordados durante el desarrollo del Trabajo Fin de Grado. En este contexto, la sostenibilidad no se entiende únicamente desde una perspectiva medioambiental, sino también desde dimensiones educativas, técnicas, sociales y profesionales.

El proyecto desarrollado consiste en la evolución de una aplicación web para la creación, validación y transformación de diagramas Entidad-Relación. Por ello, su relación con la sostenibilidad se sitúa principalmente en el ámbito de la educación, la reutilización de software, la documentación, el mantenimiento responsable y el acceso a herramientas digitales de apoyo al aprendizaje.

A continuación, se relacionan algunas de las competencias trabajadas durante el proyecto con distintos Objetivos de Desarrollo Sostenible de la Agenda 2030 de Naciones Unidas~\cite{ods-agenda2030}.

\section{Competencias de sostenibilidad adquiridas}

\textbf{ODS 4: Educación de calidad}

La aplicación desarrollada puede servir como apoyo para estudiantes que estén aprendiendo el modelo Entidad-Relación y su transformación al modelo relacional. El hecho de poder crear diagramas, validarlos y observar una posible generación de SQL facilita la práctica con ejemplos vistos en clase y permite detectar errores de forma más visual.

Además, al tratarse de una aplicación web, su uso no requiere instalar herramientas complejas. Esto reduce la barrera de entrada para el estudiante y permite centrarse en el aprendizaje del modelo. No obstante, la herramienta debe entenderse como un complemento al estudio y a la explicación del profesorado, no como un sustituto del razonamiento teórico necesario para diseñar correctamente una base de datos.

\textbf{ODS 9: Industria, innovación e infraestructura}

El proyecto se ha basado en la evolución de una herramienta software existente, incorporando nuevas funcionalidades y estabilizando su comportamiento. Esta forma de trabajo se aproxima a situaciones habituales en el ámbito profesional, donde no siempre se parte de cero, sino que es necesario comprender, mantener y ampliar sistemas ya desarrollados.

Durante el trabajo se ha seguido una organización progresiva mediante issues, milestones, pruebas y documentación. Esta metodología ha permitido abordar los cambios de forma controlada, reduciendo el riesgo de introducir modificaciones incoherentes con el resto del sistema. Desde este punto de vista, el proyecto contribuye al desarrollo de una infraestructura digital académica más mantenible y preparada para futuras ampliaciones.

\textbf{ODS 12: Producción y consumo responsables}

Uno de los aprendizajes más importantes del proyecto ha sido comprender que reutilizar software existente también implica responsabilidad. Aunque partir de una aplicación heredada puede parecer más sencillo que comenzar desde cero, en la práctica fue necesario dedicar una parte importante del trabajo a entender su estructura interna, sus decisiones de diseño y las dependencias entre el lienzo, el modelo interno, las validaciones, la transformación relacional y la generación de SQL.

Esta experiencia ha permitido valorar la reutilización como una práctica sostenible, siempre que vaya acompañada de análisis, documentación y mantenimiento. En lugar de descartar el trabajo previo, se ha aprovechado la base existente y se han introducido mejoras de forma progresiva. También se ha prestado atención a la licencia del proyecto actual y a las dependencias utilizadas, con el objetivo de respetar el trabajo de terceros y aclarar las condiciones de uso de la versión desarrollada.

\textbf{ODS 17: Alianzas para lograr los objetivos}

El ODS 17 se centra en la cooperación entre distintos actores para facilitar el desarrollo sostenible. Aunque este objetivo suele asociarse a alianzas institucionales o internacionales, también incluye aspectos relacionados con la transferencia de conocimiento, el acceso a la tecnología y la colaboración entre personas u organizaciones.

En este proyecto, la relación con el ODS 17 se plantea de forma moderada. El trabajo se apoya en software libre y en herramientas abiertas o ampliamente accesibles, como React, mxGraph, Git, GitHub, LaTeX y distintas utilidades de desarrollo y validación. Esta base tecnológica permite reutilizar conocimiento existente, reducir la dependencia de soluciones cerradas y facilitar que otros estudiantes puedan continuar el proyecto sin partir desde cero.

Además, el desarrollo ha requerido una comunicación continua con el tutor académico y una revisión progresiva de las decisiones tomadas. Esta colaboración ha sido especialmente importante porque el modelo Entidad-Relación puede admitir distintos matices de interpretación, y la aplicación debía mantenerse alineada con los materiales docentes y con el criterio seguido en la asignatura.

También se ha intentado dejar una base documentada para que futuros estudiantes o desarrolladores puedan comprender el estado del proyecto, sus funcionalidades, sus limitaciones y sus posibles líneas de continuación. En este sentido, la documentación forma parte de la sostenibilidad del trabajo, ya que facilita que el conocimiento generado no quede limitado únicamente al periodo de desarrollo del TFG.

\section{Reflexión final}

La realización de este Trabajo Fin de Grado me ha permitido trabajar competencias relacionadas con la sostenibilidad desde una perspectiva principalmente educativa y técnica. La contribución más directa del proyecto se relaciona con el ODS 4, ya que la aplicación puede servir como apoyo al aprendizaje del modelo Entidad-Relación, la validación de diagramas y la comprensión de su transformación al modelo relacional y a SQL.

También se han trabajado competencias vinculadas con el ODS 9 y el ODS 12. El proyecto no se ha desarrollado desde cero, sino a partir de una aplicación heredada que ha sido analizada, estabilizada y ampliada de forma incremental. Esta forma de trabajo refuerza la importancia de mantener software existente, reutilizar recursos técnicos, documentar las decisiones adoptadas y evitar reimplementaciones innecesarias cuando ya existe una base aprovechable.

En relación con el ODS 17, la contribución se entiende de forma prudente. El uso de software libre, la documentación del estado final del proyecto y la colaboración con el tutor académico favorecen la transferencia de conocimiento y facilitan que futuros estudiantes puedan continuar la evolución de la herramienta. No se trata de una alianza institucional amplia, sino de una contribución limitada al contexto académico del proyecto.

También es necesario reconocer algunas limitaciones. El uso de servicios en la nube para el despliegue y la validación de la aplicación puede implicar un consumo adicional de recursos frente a una ejecución exclusivamente local. Del mismo modo, una herramienta de apoyo como esta no debe sustituir el aprendizaje manual del diseño Entidad-Relación, ya que existe el riesgo de que el estudiante dependa de la aplicación sin comprender suficientemente los criterios teóricos. Por ello, la herramienta debe utilizarse como complemento docente y no como reemplazo del razonamiento necesario para diseñar correctamente una base de datos.

En conjunto, considero que el proyecto contribuye de forma modesta pero realista a la sostenibilidad educativa y técnica. Su principal aportación no está en resolver un problema medioambiental directo, sino en mejorar una herramienta académica existente, reducir barreras de uso, documentar su funcionamiento y dejar una base más mantenible para futuras ampliaciones.

